\section{Application Needs and Framework Scope}
A networked application often needs more than the location of a service. A drone operator may require a camera near an area before a deadline, while an inference requester may require a particular model, sufficient accelerator memory, and an acceptable preparation delay. A service advertisement identifies a possible source of execution, but it does not establish that the source is willing and able to undertake this task now. NDNSF separates permission to provide a service, current willingness to participate, selection for an invocation, and completion of the requested work.

These distinctions also matter within an application that usually uses one device. After choosing a drone, an operator expects subsequent camera commands to affect that drone. Substituting another available Provider would change the meaning of the command. Conversely, an image-analysis request may permit any authorized compute Provider that meets its requirements. A reusable framework must support both situations without requiring the application to implement an unrelated communication protocol for each.

The discovery descriptor contains the information needed for a participation decision. For a sensing task, this may include a coarse area, a time window, and a sensor requirement. For inference, it may include model identity, expected input size, and a deadline. Full images, detailed flight instructions, and inference prompts can be delivered after selection. The boundary follows the decision: information indispensable to assessing feasibility cannot simply be withheld until after that assessment. An application should nevertheless disclose only the necessary precision. For example, an authorized camera Provider may need an approximate search area, while an unselected candidate need not obtain the full mission route. Descriptor protection limits disclosure to authorized service Providers; it does not prevent them from learning or sharing what they legitimately decrypt.

NDNSF-UAV and NDNSF-DI expose complementary demands. The former combines device-specific operations, changing observations, live data, and physical actions. The latter combines compute selection, model preparation, large artifacts, intermediate tensors, and incremental outputs. Both require authorization and bounded request lifetimes. Their differences test whether a common invocation and named-object interface remains useful beyond one workload. Reusability is an engineering objective to be evaluated through shared mechanisms and application-specific code, rather than assumed from implementing two demonstrations.

We derive the scope by following a request from candidate discovery to execution and closure. At each stage, we identify the relationship the application needs and the failure that follows if it is not enforced.

Authorization, confidentiality, and timeliness apply across stages. Permission changes can affect an invocation throughout its lifetime and are listed separately as a cross-cutting concern.

\subsection{Research Gap and Design Focus}
NDN provides named Data retrieval and signature validation, but a Data name and a valid signature do not by themselves record or validate the service decision governing an invocation. The gap examined here is how a runtime service decision is carried into the admission and use of named Data.

Candidate discovery lets Providers report eligibility and willingness, but it does not commit an executor. A Provider's permission to offer a service authorizes that capability; it does not assign the Provider to a particular invocation or workflow role. Selection is the binding step: the initial prototype commits one Provider per invocation, while the collaboration extension commits Provider-to-role assignments and dependency edges.

For a dependent task, signature verification authenticates the producer but does not establish that the object was produced by the assigned predecessor for the active invocation or satisfies the declared dependency. The consumer must therefore check the producer, exact object name, dependency edge, and execution context.

Prior invocation and command-authorization systems address parts of this problem. NDNSF investigates a reusable protocol connecting these decisions, with authorization spanning invocation and data retrieval.

To make this gap operational, Table~\ref{tab:framework-needs} maps each invocation stage to the relationship the framework must preserve, the failure caused by its absence, and the validation planned in Chapter~\ref{chap:evaluation}.

\begin{table}[htbp]\singlespacing\centering\small
\caption{From invocation requirements to failure cases and proposed validation.}
\label{tab:framework-needs}
\begin{tabularx}{\textwidth}{p{.16\textwidth}XXp{.23\textwidth}}
\toprule Stage or concern & Required relationship & Failure if missing & Proposed checks\\\midrule
Candidate discovery & Eligible candidates report willingness and feasibility; descriptor access is restricted & Ineligible candidate selected or descriptor disclosed to unauthorized readers & ACK eligibility and descriptor access\\\midrule
Selection and start & The intended Provider starts only a permitted, active invocation & Wrong device, unauthorized work, or repeated start & Permission, Selection, and replay checks\\\midrule
Input and intermediate Data & Correct producer, object or stream, and permitted dependency & Wrong, stale, or incomplete input advances execution & Object/stream validation; role, edge, and cache tests\\\midrule
Closure or replacement & Results match the active operation; closed work cannot advance its replacement & Late results revive closed work or complete the wrong operation & Deadline, cancellation, and replacement tests\\\midrule
Permission changes & Permission updates take effect after status validation and required key updates & A grant is unusable or withdrawn permission remains accepted & Grant/withdrawal propagation and rejection delay\\\bottomrule
\end{tabularx}
\end{table}

This table is a requirements-to-evidence map; it does not claim that all proposed collaboration checks have already been implemented or validated.

The evaluation asks whether its checks enforce the intended relationships and whether the resulting cost is justified for the studied workloads.

For each workflow in Chapters~\ref{chap:uav} and~\ref{chap:di}, validation will identify the applicable rows, test their failure cases through the corresponding framework paths, and record the outcomes. Single-Provider workflows do not require collaboration-role checks.

Chapter~\ref{chap:evaluation} specifies the proposed tests. Coverage is limited to the declared workflows; the exclusions below define the remaining scope boundaries.
